\documentclass[12pt]{article}

% pacotes obrigatorios e formatacao
\usepackage[utf8]{inputenc}
\usepackage[T1]{fontenc}
\usepackage[brazil]{babel}
\usepackage[margin=2cm]{geometry} % margem minima de 2cm exigida pela metodologia
\usepackage{times} % fonte times new roman
\usepackage{graphicx}
\usepackage{booktabs}
\usepackage{xcolor}
\usepackage{listings}
\usepackage{fancyhdr}
\usepackage{hyperref}

% ---------------
% Resultados da enumeração. Use o comando abaixo para inserir screenshots:
% \begin{figure}[h!]
%     \centering
%     \includegraphics[width=0.8\textwidth]{screenshot_nmap.png}
%     \caption{Legenda descritiva da imagem inserida imediatamente após o texto.}
%     \label{fig:nmap}
% \end{figure}

% \begin{lstlisting}[caption={Comando de exploração utilizado}]
% # Exemplo de bloco de código com fundo destacado
% nmap -sV -sC -oN scan.txt 10.10.10.10
% \end{lstlisting}
% ---------------

\definecolor{citapurple}{HTML}{6A1B9A} % tom roxo escuro
% configuracao de links (hierarquia e cores)
\hypersetup{
    colorlinks=true,
    linkcolor=black,       % cor para sumario, figuras (\ref) e secoes
    urlcolor=blue,         % cor para links de internet (\url, \href)
    citecolor=citapurple,  % cor para citacoes bibliograficas (\cite) em roxo
    pdftitle={WriteUp - CTF},
}


% configuracao dos blocos de codigo (monoespacado com fundo destacado)
\definecolor{codebg}{gray}{0.95}
\lstset{
    backgroundcolor=\color{codebg},
    basicstyle=\ttfamily\small,
    breaklines=true,
    frame=single,
    frameround=tttt,
    framesep=5pt,
    rulecolor=\color{gray!30},
    showstringspaces=false,
    keywordstyle=\color{blue},
    captionpos=b,
    inputencoding=utf8,
    extendedchars=true
}

% configuracao de cabecalho e rodape
\pagestyle{fancy}
\fancyhf{}
\lhead{\small \textit{Nome da Maquina - Plataforma}}
\rhead{\small \textit{Autor: nome (nickname)}}
\rfoot{\thepage} % numeracao de paginas obrigatoria

\begin{document}

% 2.1. capa
\begin{titlepage}
    \centering
    \vspace*{2cm}
    
    {\Huge \textbf{Nome da Máquina: Write Up}}\\[0.5cm]
    {\large \textbf{Plataforma:} HackTheBox}\\[1.5cm]
    
    \vfill
    
    % espaco dedicado para logotipo (plataforma, unbreakable ou da máquina)
    \begin{figure}[h!]
        \centering
        % substitua 'example-image' pelo arquivo da logo na pasta do projeto
        \includegraphics[width=0.32\textwidth]{unbreakableLogo_3.png} 
    \end{figure}
    
    \vfill
    
    \begin{table}[h!]
        \centering
        \small
        \renewcommand{\arraystretch}{1.3}
        \begin{tabular}{ll}
            \textbf{Máquina / Desafio:} & Nome da Máquina \\
            \textbf{Dificuldade:}       & Fácil / Média / Difícil / Insane \\
            \textbf{Autor:}             & Lurian Lima (luricl) \\
            \textbf{Matrícula:}         & 222005395 \\
        \end{tabular}
    \end{table}
    
    \vspace*{1.5cm}
    {\large \today}
\end{titlepage}

\newpage

% 2.2. sumario
% indice gerado automaticamente com numeracao de paginas
\tableofcontents
\newpage

% 2.3. visao geral
% paragrafo descrevendo o vetor principal de ataque sem spoiler detalhado.
% deve responder: o que e o desafio, qual a superficie explorada e o caminho geral ate o objetivo.
\section{Visão Geral}
A máquina apresenta um cenário focado na exploração de falhas em aplicações web e configurações inadequadas no sistema operacional. O acesso inicial foi obtido através da identificação de vulnerabilidades na autenticação da aplicação, permitindo a execução remota de código. Em seguida, a escalada de privilégios explorou permissões locais mal configuradas para obter controle total do servidor como \texttt{root}.


% 2.4. reconhecimento
% resultados da enumeracao externa: portas abertas, servicos e versoes relevantes identificadas.
% incluir comandos executados em blocos de codigo e outputs truncados pertinentes.
% regra: nao incluir informacoes obtidas apos o foothold inicial.
\section{Reconhecimento}

Iniciou-se a varredura de portas com o Nmap para mapear os servicos ativos no alvo, seguindo as diretrizes descritas em \cite{ref1}:

\begin{lstlisting}[caption={Varredura de portas com Nmap}, label={lst:nmap}]
nmap -sV <TARGET_IP>
\end{lstlisting}

O resultado da varredura pode ser visualizado na Figura \ref{fig:nmap}, evidenciando os servicos encontrados.

\begin{figure}[h!]
    \centering
    % substitua 'example-image' pelo screenshot real
    \includegraphics[width=0.45\textwidth]{example-image}
    \caption{Servicos identificados no alvo via Nmap.}
    \label{fig:nmap}
\end{figure}

% 2.5. exploracao
% descricao tecnica detalhada do vetor de entrada e raciocinio que levou a vulnerabilidade.
% incluir scripts, exploits ou payloads utilizados (comentados) e justificativa dos passos.
% prints obrigatorios evidenciando o acesso inicial obtido (com placeholder para dados sensiveis).
\section{Exploração}


% 2.6. movimentacao lateral (se aplicavel)
% aplicavel quando ha multiplos usuarios no sistema ou necessidade de pivoting entre servicos.
% detalhar a escalada horizontal com comandos, raciocinio e evidencias necessarias.
\section{Movimentação Lateral}


% 2.7. escalada de privilegios
% descrever o vetor de privesc vertical, o processo de enumeracao local que o revelou
% (ex: linpeas, sudo -l, capabilities, binarios suid) e a exploracao realizada.
% incluir o comando final e captura que comprove o acesso root/system e as flags com placeholder.
\section{Escalada de Privilégios}


% 2.8. conclusao
% resumo tecnico de 3 a 5 linhas sintetizando a cadeia de exploracao completa,
% do reconhecimento a captura final da flag.
\section{Conclusão}

\section{Referências}

% 2.9. referencias
% espaco destinado as fontes utilizadas para a exploracao:
% links dos repositorios de exploits publicos, documentacoes e cves correspondentes.
\begin{thebibliography}{9}
    \bibitem{ref1}
    \textbf{OWASP / CVE / Nome da Vulnerabilidade}: Breve descrição técnica do vetor abordado, documentação ou payload de referência. \\
    Disponível em: \url{https://exemplo.com/artigo-ou-documentacao}
    
    \bibitem{ref2}
    \textbf{Repositório / Autor do Exploit}: Detalhes do exploit público, script ou ferramenta utilizada durante a exploração. \\
    Disponível em: \url{https://github.com/exemplo/repositorio-exploit}

\end{thebibliography}

\end{document}